\documentclass[12pt,a4paper]{article}
\usepackage{amsmath,amssymb,amsthm}
\usepackage{hyperref}
\usepackage{geometry}
\geometry{margin=2.5cm}
\usepackage{enumitem}
\usepackage{booktabs}

\newtheorem{theorem}{Theorem}[section]
\newtheorem{lemma}[theorem]{Lemma}
\newtheorem{corollary}[theorem]{Corollary}
\newtheorem{proposition}[theorem]{Proposition}
\theoremstyle{definition}
\newtheorem{definition}[theorem]{Definition}
\theoremstyle{remark}
\newtheorem{remark}[theorem]{Remark}

\title{Modal Logic Grounded in Cost Geometry:\\
Necessity, Possibility, and Actuality from $J(x)$}

\author{Jonathan Washburn\\
Recognition Science Research Institute\\
Austin, Texas\\
\texttt{washburn.jonathan@gmail.com}}

\date{March 2026}

\begin{document}
\maketitle

\begin{abstract}
We ground the fundamental concepts of modal logic---necessity,
possibility, impossibility, and actuality---in the cost functional
$J(x) = \tfrac{1}{2}(x + x^{-1}) - 1$, the unique normalized solution
to the Recognition Composition Law.  A configuration is \emph{necessary}
if and only if it is the unique $J$-minimizer ($x = 1$, the sole
zero-defect state).  A configuration is \emph{possible} if and only if
it has positive ratio ($x > 0$, ensuring finite $J$-cost).  A
configuration is \emph{impossible} if $x \leq 0$ (violating ledger
positivity).  Actuality is identified with the $J$-minimum---a form of
necessitarianism grounded in cost geometry.  We prove that the core S5
structural content holds as theorems: necessity implies actuality implies
possibility; the necessary configuration is unique; and every possible
configuration is accessible to the necessary one via cost-decreasing
paths.  We further define propositional modal operators $\Box$ and
$\Diamond$ over configuration space and prove the distribution axiom~K,
modal duality, and a weak form of the truth axiom~T.  Modal distance
from the actual world is measured by defect: $d(x) = J(x) \geq 0$, with
$d(x) = 0$ if and only if $x = 1$.  This provides modal metaphysics
with physical content: necessity is not ``truth in all possible worlds''
but unique cost-geometric forcing; possibility is not a Lewisian concrete
world but finite cost; and actuality is not indexical but the
$J$-minimum.  The foundational mathematical results (uniqueness of the
$J$-minimizer, non-negativity of $J$, AM--GM inequality) are
rigorously proved within the RS framework.
\end{abstract}

%% ============================================================
\section{Introduction}
\label{sec:intro}
%% ============================================================

Modal logic---the logic of necessity ($\Box$) and possibility
($\Diamond$)---is central to metaphysics, epistemology, and the
philosophy of language.  Yet its foundational concepts remain contested.
What does it mean for something to be ``necessary''?  What are
``possible worlds''?  Is the actual world distinguished structurally, or
is its status merely indexical?

Lewis~\cite{Lewis1986} proposed modal realism: possible worlds are
concrete entities as real as the actual world, and actuality is indexical
(``the world I happen to inhabit'').  Kripke~\cite{Kripke1980} treated
possible worlds as abstract stipulations, grounding necessity in rigid
designation.  Plantinga~\cite{Plantinga1974} identified necessity with
maximal states of affairs.  Williamson~\cite{Williamson2013} developed a
systematic case for treating modal logic as first-order metaphysics,
arguing that necessitism---the thesis that everything necessarily
exists---provides the best overall theory.
Stalnaker~\cite{Stalnaker2012} offered a moderate realism about
propositions as sets of possible worlds while resisting Lewis's
ontological inflation.  Each approach has virtues, but none derives modal
concepts from a deeper mathematical structure with independently
motivated physical content.

Recognition Science~\cite{RS_Axioms} provides such a derivation.
The theory begins with a single functional equation---the Recognition
Composition Law (RCL)---whose unique normalized solution is the cost
functional
$J(x) = \tfrac{1}{2}(x + x^{-1}) - 1$ on $\mathbb{R}_{>0}$.  From
this functional alone, we derive precise mathematical content for each
modal concept:

\begin{itemize}[nosep]
  \item \textbf{Necessity} $=$ unique $J$-minimizer ($x = 1$).
  \item \textbf{Possibility} $=$ positive ratio ($x > 0$, finite cost).
  \item \textbf{Impossibility} $=$ non-positive ratio ($x \leq 0$,
  undefined cost).
  \item \textbf{Actuality} $=$ the $J$-minimum itself, not an indexical label.
  \item \textbf{Modal distance} $=$ defect $J(x)$, a continuous
  quantitative measure.
\end{itemize}

The resulting modal structure satisfies the S5 axioms as \emph{theorems}
rather than postulates, and replaces the metaphysical debate about
possible worlds with cost geometry.  The framework entails a form of
necessitarianism---the actual world is the only necessary one---that
echoes Spinoza~\cite{Spinoza1677} and resolves the Leibniz question.

The paper is structured as follows.  Section~\ref{sec:cost} derives the
cost functional.  Section~\ref{sec:defs} introduces the RS modal
definitions.
Sections~\ref{sec:necessity}--\ref{sec:actuality} prove the structural
theorems.  Section~\ref{sec:s5} demonstrates the S5 axioms.
Section~\ref{sec:distance} develops quantitative modality.
Section~\ref{sec:operators} defines propositional modal operators.
Section~\ref{sec:discuss} compares the framework with standard
approaches and draws philosophical consequences.

%% ============================================================
\section{The Cost Functional}
\label{sec:cost}
%% ============================================================

\subsection{Axioms and Uniqueness}

The cost functional is the unique solution to three axioms:

\noindent\textbf{Axiom A1 (Normalization):} $J(1) = 0$ (the identity
configuration has zero cost).\\
\textbf{Axiom A2 (Recognition Composition Law, RCL):} For all
$x, y > 0$:
\begin{equation}
\label{eq:RCL}
  J(xy) + J(x/y) = 2J(x)J(y) + 2J(x) + 2J(y).
\end{equation}
This equation expresses that the cost of comparing a composite ($xy$)
and its reciprocal complement ($x/y$) decomposes multiplicatively and
additively in terms of the component costs.\\
\textbf{Axiom A3 (Calibration):} $J''_{\log}(0) = 1$, where
$J_{\log}(t) := J(e^t)$.

\begin{theorem}[Cost Uniqueness]
\label{thm:cost-unique}
The continuous solution to \textup{(A1)--(A3)} is unique:
$J(x) = \tfrac{1}{2}(x + x^{-1}) - 1$.
\end{theorem}

\begin{proof}[Proof sketch]
Setting $f(t) = J_{\log}(t) + 1$ and rewriting \textup{(A2)} gives
the d'Alembert functional equation $f(s+t) + f(s-t) = 2f(s)f(t)$.
By the Acz\'el classification theorem~\cite{Aczel1966}, the unique
continuous solution with $f(0) = 1$ and $f''(0) = 1$ is
$f(t) = \cosh t$.  Inverting, $J(x) = f(\ln x) - 1 = \cosh(\ln x) - 1
= \tfrac{1}{2}(x + x^{-1}) - 1$.
\end{proof}

\subsection{Verification and Key Properties}

Subject to (A1)--(A3), the unique $C^2$ solution is:

\begin{definition}[Cost Functional and Defect]
\label{def:Jx}
For $x \in \mathbb{R}_{>0}$:
\begin{equation}
\label{eq:J}
  J(x) = \frac{1}{2}\bigl(x + x^{-1}\bigr) - 1.
\end{equation}
The \emph{defect} of $x$ is $\mathrm{defect}(x) := J(x)$.
\end{definition}

\noindent
That this $J$ satisfies the RCL can be verified by direct computation.
The left side:
\begin{equation}
  J(xy) + J(x/y) = \tfrac{1}{2}\bigl(xy + (xy)^{-1} + xy^{-1}
  + x^{-1}y\bigr) - 2.
\end{equation}
The right side:
\begin{equation}
  2J(x)J(y) + 2J(x) + 2J(y)
  = \tfrac{1}{2}(x + x^{-1})(y + y^{-1}) - 2,
\end{equation}
where the intermediate terms cancel.  Expanding the product on the right
gives $\tfrac{1}{2}(xy + xy^{-1} + x^{-1}y + x^{-1}y^{-1}) - 2$,
confirming equality.

\subsection{Key Properties}

\begin{lemma}[Basic Properties of $J$]
\label{lem:J-props}
\begin{enumerate}[label=(\roman*)]
  \item $J(x) \geq 0$ for all $x > 0$, with equality if and only if
  $x = 1$.
  \item $J(x) = J(x^{-1})$ for all $x > 0$ (reciprocal symmetry).
  \item $J$ is strictly convex on $(0, \infty)$, with
  $J''(x) = x^{-3} > 0$.
  \item $J(x) \to \infty$ as $x \to 0^+$ or $x \to \infty$.
\end{enumerate}
\end{lemma}

\begin{proof}
(i) By AM--GM, $x + x^{-1} \geq 2\sqrt{x \cdot x^{-1}} = 2$ for
$x > 0$, with equality iff $x = x^{-1}$, i.e., $x = 1$.  Hence
$J(x) \geq 0$.

(ii) $J(x^{-1}) = \tfrac{1}{2}(x^{-1} + x) - 1 = J(x)$.

(iii) $J'(x) = \tfrac{1}{2}(1 - x^{-2})$ and
$J''(x) = x^{-3} > 0$ for $x > 0$.

(iv) As $x \to 0^+$, $x^{-1} \to \infty$, so $J(x) \to \infty$.
As $x \to \infty$, $\tfrac{1}{2}x \to \infty$, so $J(x) \to \infty$.
\end{proof}

%% ============================================================
\section{RS Modal Definitions}
\label{sec:defs}
%% ============================================================

We define the four fundamental modal concepts in terms of $J$.  The
geometric motivation is that $J$ defines a cost landscape on
$\mathbb{R}_{>0}$ with a unique global minimum at $x = 1$.  Modal
concepts correspond to positions on this landscape.

\begin{definition}[RS Necessity]
\label{def:necessary}
A configuration $x \in \mathbb{R}$ is \emph{RS-necessary} if and only if
\begin{equation}
  x > 0 \quad\text{and}\quad J(x) = 0.
\end{equation}
\end{definition}

\noindent
A configuration is necessary if it incurs zero cost.  It is not
``forced'' by fiat but by the structure of $J$---any deviation from
$x = 1$ strictly increases cost.

\begin{definition}[RS Possibility]
\label{def:possible}
A configuration $x$ is \emph{RS-possible} if and only if $x > 0$.
\end{definition}

\noindent
A configuration is possible if it has a well-defined, finite cost.  The
positivity requirement ensures that $J(x)$ is defined and real.
Configurations with $x \leq 0$ lie outside the domain of $J$ entirely.

\begin{definition}[RS Actuality]
\label{def:actual}
A configuration $x$ is \emph{RS-actual} if and only if it is
RS-necessary.
\end{definition}

\noindent
This identification---actuality \emph{is} necessity---is the central
metaphysical claim.  It is not a tautology but a substantive thesis:
the actual world is not one possible world among many that happens to be
``ours'' (Lewis), nor a stipulated entity (Kripke), but the unique
cost-geometric minimum.  The philosophical significance of this
identification is discussed in Section~\ref{sec:discuss}.

\begin{definition}[RS Impossibility]
\label{def:impossible}
A configuration $x$ is \emph{RS-impossible} if and only if $x \leq 0$.
\end{definition}

\begin{definition}[Modal Distance]
\label{def:distance}
The \emph{modal distance} from $x$ to the actual world is
\begin{equation}
  d(x) := J(x) = \frac{1}{2}\bigl(x + x^{-1}\bigr) - 1.
\end{equation}
\end{definition}

%% ============================================================
\section{Necessity: The Unique $J$-Minimizer}
\label{sec:necessity}
%% ============================================================

\begin{theorem}[Unique Necessity]
\label{thm:unique-nec}
There exists exactly one RS-necessary configuration: $x = 1$.
\end{theorem}

\begin{proof}
\textit{Existence.}  $1 > 0$ and $J(1) = \tfrac{1}{2}(1 + 1) - 1 = 0$.

\textit{Uniqueness.}  If $y > 0$ and $J(y) = 0$, then
$\tfrac{1}{2}(y + y^{-1}) = 1$, so $y + y^{-1} = 2$, so
$y^2 - 2y + 1 = 0$, so $(y - 1)^2 = 0$, so $y = 1$.
\end{proof}

\begin{corollary}[Necessity is Identity]
\label{cor:nec-one}
For all $x \in \mathbb{R}$: $x$ is RS-necessary if and only if $x = 1$.
\end{corollary}

\begin{proof}
Forward: Theorem~\ref{thm:unique-nec}.  Backward: $J(1) = 0$ and
$1 > 0$.
\end{proof}

\begin{remark}[Structural Forcing]
\label{rem:forcing}
The uniqueness of the necessary configuration follows from the strict
convexity of $J$ (Lemma~\ref{lem:J-props}(iii)).  A strictly convex
function on a convex domain has at most one global minimum.  Since $J$
achieves its minimum at $x = 1$ and is strictly convex on
$(0, \infty)$, no other configuration can have zero defect.  Necessity
is \emph{structurally forced} by the geometry of the cost landscape, not
postulated.
\end{remark}

%% ============================================================
\section{Possibility: Positive Ratio and Finite Cost}
\label{sec:possibility}
%% ============================================================

\begin{theorem}[Possibility is Broad]
\label{thm:possible}
For all $x > 0$: $x$ is RS-possible.  Conversely, every RS-possible
$x$ satisfies $x > 0$.
\end{theorem}

\begin{proof}
Immediate from Definition~\ref{def:possible}.
\end{proof}

\begin{theorem}[Finite Cost for Possible Configurations]
\label{thm:finite-cost}
Every RS-possible configuration has non-negative, finite $J$-cost.
\end{theorem}

\begin{proof}
For $x > 0$, $J(x) = \tfrac{1}{2}(x + x^{-1}) - 1$ is a finite real
number, and $J(x) \geq 0$ by Lemma~\ref{lem:J-props}(i).
\end{proof}

\begin{theorem}[Possibility is Not Singleton]
\label{thm:not-singleton}
There exist distinct RS-possible configurations.
\end{theorem}

\begin{proof}
$x = 1$ and $x = 2$ are both positive, hence both RS-possible, and
$1 \neq 2$.
\end{proof}

%% ============================================================
\section{Impossibility: The Positivity Boundary}
\label{sec:impossible}
%% ============================================================

\begin{theorem}[Impossibility Equivalence]
\label{thm:impossible}
$x$ is RS-impossible if and only if $x$ is not RS-possible.
\end{theorem}

\begin{proof}
RS-impossible means $x \leq 0$.  Not RS-possible means $\neg(x > 0)$,
which is $x \leq 0$.
\end{proof}

\begin{theorem}[No Overlap]
\label{thm:no-overlap}
No configuration is both RS-possible and RS-impossible.
\end{theorem}

\begin{proof}
$x > 0$ and $x \leq 0$ are contradictory.
\end{proof}

\begin{remark}[The Positivity Boundary]
The boundary $x = 0$ represents a hard limit: as $x \to 0^+$,
$J(x) \to \infty$ (Lemma~\ref{lem:J-props}(iv)), so the boundary is
not merely forbidden but infinitely expensive to approach.
Impossibility thus has quantitative character: the ``closer'' a
configuration is to the boundary (in the sense of small $x$), the more
costly it is.
\end{remark}

%% ============================================================
\section{Actuality: The $J$-Minimum}
\label{sec:actuality}
%% ============================================================

\begin{theorem}[Actuality is Unique]
\label{thm:actual-unique}
$x = 1$ is the unique RS-actual configuration.
\end{theorem}

\begin{proof}
RS-actual $\Leftrightarrow$ RS-necessary $\Leftrightarrow$ $x = 1$
(Corollary~\ref{cor:nec-one}).
\end{proof}

\begin{theorem}[The Merely Possible]
\label{thm:merely-possible}
For $x > 0$ with $x \neq 1$: $x$ is RS-possible but not RS-actual.
\end{theorem}

\begin{proof}
$x > 0$ ensures RS-possible.  $x \neq 1$ implies $J(x) > 0$
(Lemma~\ref{lem:J-props}(i)), so $x$ is not RS-necessary, hence not
RS-actual.
\end{proof}

\begin{remark}[Tripartite Classification]
The space of configurations is partitioned into three classes:
\begin{enumerate}[nosep,label=(\alph*)]
  \item The unique \emph{actual/necessary} configuration: $x = 1$,
  $J(x) = 0$.
  \item The \emph{merely possible} configurations: $x > 0$, $x \neq 1$,
  $J(x) > 0$.
  \item The \emph{impossible} configurations: $x \leq 0$.
\end{enumerate}
This classification emerges from the cost landscape alone.
\end{remark}

%% ============================================================
\section{S5 Modal Axioms as Theorems}
\label{sec:s5}
%% ============================================================

The standard S5 system of modal logic consists of the axioms K
(distribution), T (truth), and 5 (positive introspection), or
equivalently K, T, and B (symmetry).  In standard Kripke semantics, S5
corresponds to an accessibility relation that is an equivalence relation
(reflexive, symmetric, transitive).  We show that the RS cost geometry
satisfies the structural content of S5.

\subsection{The Implication Chain}

\begin{theorem}[Axiom T: Necessity $\Rightarrow$ Actuality]
\label{thm:T}
For all $x$: if $x$ is RS-necessary, then $x$ is RS-actual.
\end{theorem}

\begin{proof}
RS-actual is defined as RS-necessary (Definition~\ref{def:actual}).
\end{proof}

\begin{remark}
The identification actuality $=$ necessity (Definition~\ref{def:actual})
makes Theorem~\ref{thm:T} definitionally immediate.  The substantive
content is in the definition itself: the claim that the unique
$J$-minimum plays the role of the ``actual world'' gives modal Axiom~T a
geometric rather than a stipulative basis.
\end{remark}

\begin{theorem}[Axiom D: Actuality $\Rightarrow$ Possibility]
\label{thm:D}
For all $x$: if $x$ is RS-actual, then $x$ is RS-possible.
\end{theorem}

\begin{proof}
RS-actual implies RS-necessary, which requires $x > 0$ by
Definition~\ref{def:necessary}.  Hence $x$ is RS-possible by
Definition~\ref{def:possible}.
\end{proof}

\begin{corollary}[Full Implication Chain]
\label{cor:chain}
For all $x$:
\[
  \text{RS-necessary}(x) \;\Longrightarrow\; \text{RS-actual}(x)
  \;\Longrightarrow\; \text{RS-possible}(x).
\]
\end{corollary}

\subsection{Accessibility}

The S5 axiom B (symmetry of accessibility) is captured in the RS
framework by the following theorem, which is strictly stronger than B:
not only is every possible configuration accessible to the necessary
one, but the access path is cost-monotone.

\begin{theorem}[Cost-Monotone Accessibility]
\label{thm:B}
For every RS-possible $x$, there exists a continuous path
$\gamma : [0, 1] \to (0, \infty)$ with $\gamma(0) = x$,
$\gamma(1) = 1$, and $J(\gamma(t))$ non-increasing.
\end{theorem}

\begin{proof}
Define $\gamma(t) = x^{1-t}$ for $t \in [0, 1]$.  Then
$\gamma(0) = x$, $\gamma(1) = x^0 = 1$, and $\gamma(t) = x^{1-t} > 0$
for all $t$ (since $x > 0$).

Let $s = 1 - t$ and $u(s) = J(x^s) = \tfrac{1}{2}(x^s + x^{-s}) - 1$.
Then
\[
  u'(s) = \tfrac{1}{2}\ln(x)\bigl(x^s - x^{-s}\bigr).
\]
For $x > 1$: $\ln x > 0$ and $x^s > x^{-s}$ when $s > 0$, so
$u'(s) > 0$.  Since $s = 1 - t$, we have
$\frac{d}{dt} J(\gamma(t)) = -u'(s) < 0$ for $t < 1$.

For $0 < x < 1$: $\ln x < 0$ and $x^s < x^{-s}$ when $s > 0$, so
$u'(s) = (\text{negative})(\text{negative}) > 0$.  Again
$\frac{d}{dt} J(\gamma(t)) < 0$.

For $x = 1$: $\gamma(t) \equiv 1$ and $J \equiv 0$.

In all cases, $J(\gamma(t))$ is non-increasing from $J(x)$ to $J(1) = 0$.
\end{proof}

\begin{remark}[Structural S5]
In the RS framework, the accessibility structure is richer than an
equivalence relation: it is a continuous cost landscape with a unique
global minimum.  Theorem~\ref{thm:B} provides not merely that every
possible configuration \emph{can reach} the necessary one, but that
there is a \emph{directed}, cost-decreasing path.  This gives the
accessibility relation a natural orientation absent from standard Kripke
frames.
\end{remark}

%% ============================================================
\section{Modal Distance: Quantitative Modality}
\label{sec:distance}
%% ============================================================

\begin{theorem}[Modal Distance Properties]
\label{thm:distance}
\begin{enumerate}[label=(\roman*)]
  \item $d(x) \geq 0$ for all $x > 0$ (non-negativity).
  \item $d(x) = 0$ if and only if $x = 1$, i.e., $x$ is RS-actual
  (identity of indiscernibles).
  \item $d(x) = d(x^{-1})$ for all $x > 0$ (reciprocal symmetry).
  \item $d(1) = 0$ (zero self-distance).
\end{enumerate}
\end{theorem}

\begin{proof}
(i)--(ii) follow from Lemma~\ref{lem:J-props}(i).
(iii) follows from Lemma~\ref{lem:J-props}(ii).
(iv) is $J(1) = 0$.
\end{proof}

\begin{remark}[Quantitative Modality]
Standard modal logic is qualitative: a proposition is necessary,
possible, or impossible, with no intermediate grades.  Modal distance
replaces this trichotomy with a continuous spectrum.  A configuration
with $d(x) = 0.01$ is ``almost necessary''; one with $d(x) = 100$ is
``very far from the actual world.''  This quantitative structure has no
analog in Lewis, Kripke, Plantinga, or Williamson.

The reciprocal symmetry $d(x) = d(x^{-1})$ means that ``doubling'' and
``halving'' are equidistant from actuality.  A world where some ratio is
$2$ is exactly as far from actuality as a world where that ratio is
$\tfrac{1}{2}$.  This is a non-obvious structural feature inherited from
the cost functional's reciprocal invariance.
\end{remark}

%% ============================================================
\section{Propositional Modal Operators}
\label{sec:operators}
%% ============================================================

The definitions in Sections~\ref{sec:defs}--\ref{sec:actuality}
concern \emph{configurations} (real numbers).  We now lift these to
\emph{propositional} modal operators over a configuration space, making
contact with standard modal logic syntax.

\subsection{Configuration Space and Possibility Set}

\begin{definition}[Configuration]
\label{def:config}
A \emph{configuration} is a triple $c = (x, t, B)$ where $x > 0$ is
the ratio value, $t \in \mathbb{N}$ is the tick count, and
$|\!\ln x| \leq B$ for a fixed bound $B \leq 16$.  We write
$\mathcal{C}$ for the set of all configurations.
\end{definition}

\begin{remark}[The bound $B \leq 16$]
The restriction $B \leq 16$ in Definition~\ref{def:config} is required
by Lemma~\ref{lem:identity-possible}: for the identity $A(c)$ to be
accessible from $c$, one needs $|\!\ln x| \leq 16$ (see proof below).
Physically, $|\!\ln x| \leq 16$ means $e^{-16} \leq x \leq e^{16}$
(i.e., ratios within a factor of $\approx 8.9 \times 10^6$ of unity),
which is a generous bound covering all RS physical ratios.  Configurations
outside this range require multi-epoch transitions and are not excluded
from RS physics, but they fall outside the single-epoch modal framework
here.
\end{remark}

\begin{definition}[Stasis and Transition Costs]
\label{def:costs}
For configurations with values $x, y > 0$:
\begin{align}
  J_{\mathrm{stasis}}(x) &= 8 \cdot J(x), \\
  J_{\mathrm{trans}}(x, y) &= |\!\ln(y/x)| \cdot
  \tfrac{1}{2}\bigl(J(x) + J(y)\bigr).
\end{align}
The stasis cost is the $J$-cost of remaining at $x$ for one eight-tick
epoch.  The transition cost measures the expense of changing from $x$
to $y$, weighted by the average defect.
\end{definition}

\begin{definition}[Possibility Set]
\label{def:poss-set}
For a configuration $c = (x, t, B)$, the \emph{possibility set}
$\mathrm{Poss}(c)$ consists of all configurations $c' = (y, t{+}8, B')$
such that transitioning does not exceed the stasis cost:
\begin{equation}
  J_{\mathrm{trans}}(x, y) + J(y) \leq J_{\mathrm{stasis}}(x).
\end{equation}
\end{definition}

\begin{definition}[Actualization]
\label{def:actualize}
The \emph{actualization} of $c = (x, t, B)$ is
$A(c) = (1, t{+}8, 0)$---the identity configuration at the next epoch.
\end{definition}

\begin{lemma}[The Identity is Always Possible]
\label{lem:identity-possible}
For every $c \in \mathcal{C}$, $A(c) \in \mathrm{Poss}(c)$.
\end{lemma}

\begin{proof}
$J(1) = 0$, so the constraint becomes
$J_{\mathrm{trans}}(x, 1) + 0 \leq 8\,J(x)$, i.e.,
$|\!\ln x| \cdot \tfrac{J(x)}{2} \leq 8\,J(x)$.  If $J(x) = 0$ then
$x = 1$ and both sides vanish.  If $J(x) > 0$, dividing gives
$|\!\ln x| \leq 16$, which holds by the bound constraint on
configurations.
\end{proof}

\subsection{Modal Operators}

\begin{definition}[Box and Diamond]
\label{def:box-diamond}
For a configuration property
$p : \mathcal{C} \to \{\mathrm{T}, \mathrm{F}\}$:
\begin{align}
  (\Box p)(c) &\;\Longleftrightarrow\;
  \forall\, c' \in \mathrm{Poss}(c),\; p(c'), \\
  (\Diamond p)(c) &\;\Longleftrightarrow\;
  \exists\, c' \in \mathrm{Poss}(c),\; p(c').
\end{align}
\end{definition}

\subsection{Key Theorems}

\begin{proposition}[Modal Duality]
\label{prop:duality}
$(\Box p)(c) \;\Longleftrightarrow\; \neg\,(\Diamond\,(\neg p))(c)$.
\end{proposition}

\begin{proof}
$(\Box p)(c)$ means $\forall c' \in \mathrm{Poss}(c),\, p(c')$, which
is $\neg\,\exists c' \in \mathrm{Poss}(c),\, \neg p(c')$, which is
$\neg\,(\Diamond\,(\neg p))(c)$.
\end{proof}

\begin{proposition}[Axiom K: Distribution]
\label{prop:K}
If\/ $(\Box(p \Rightarrow q))(c)$ and $(\Box p)(c)$, then
$(\Box q)(c)$.
\end{proposition}

\begin{proof}
Let $c' \in \mathrm{Poss}(c)$.  By hypothesis,
$p(c') \Rightarrow q(c')$ and $p(c')$.  Hence $q(c')$.  Since $c'$
was arbitrary, $(\Box q)(c)$.
\end{proof}

\begin{proposition}[Weak Axiom T]
\label{prop:T-weak}
If $(\Box p)(c)$, then $p(A(c))$.
\end{proposition}

\begin{proof}
By Lemma~\ref{lem:identity-possible}, $A(c) \in \mathrm{Poss}(c)$.
Since $(\Box p)(c)$ asserts $p(c')$ for all $c' \in \mathrm{Poss}(c)$,
we obtain $p(A(c))$.
\end{proof}

\begin{remark}[Temporal Asymmetry and Weak T]
This is a ``weak'' form of Axiom T because it asserts that necessity
implies truth at the \emph{actualized} configuration $A(c)$, not at
$c$ itself.  The standard Axiom T ($\Box p \to p$) would require $c$ to
be in its own possibility set, which the time-advance structure
($t \to t{+}8$) prevents: $c$ has tick $t$, while all elements of
$\mathrm{Poss}(c)$ have tick $t{+}8$.  This temporal asymmetry is a
feature rather than a defect: it reflects the arrow of time.  The
identity configuration at the \emph{next} tick plays the role of the
``actual world'' in the standard T axiom.
\end{remark}

\begin{proposition}[Actualization Decreases Cost]
\label{prop:cost-decrease}
If $c = (x, t, B)$ with $x \neq 1$, then
$J(A(c).\,x) < J(c.\,x)$, i.e., $0 < J(x)$.
\end{proposition}

\begin{proof}
$A(c).\,x = 1$ and $J(1) = 0 < J(x)$ for $x \neq 1$
(Lemma~\ref{lem:J-props}(i)).
\end{proof}

%% ============================================================
\section{Discussion}
\label{sec:discuss}
%% ============================================================

\subsection{Comparison with Standard Frameworks}

\begin{center}
\renewcommand{\arraystretch}{1.3}
\begin{tabular}{@{}lcccc@{}}
\toprule
& \textbf{Lewis} & \textbf{Kripke} & \textbf{Williamson} & \textbf{RS} \\
\midrule
Possible worlds & Concrete & Abstract & Abstract & Finite-cost configs \\
Necessity & All-worlds truth & Rigid desig. & Necessitism & Unique $J$-min \\
Actuality & Indexical & Privileged & Distinguished & $J$-minimum \\
S5 axioms & Postulated & Postulated & Argued & Proved \\
Ontological posit & Many worlds & Possible worlds & Necessary beings & None extra \\
Modal distance & None & None & None & $d(x) = J(x)$ \\
\bottomrule
\end{tabular}
\end{center}

RS has several advantages over each standard framework:
\begin{itemize}
  \item \textbf{Over Lewis.}  No ontological inflation.  RS does not
  posit concrete possible worlds---merely configurations with positive
  cost.  Lewis's ``plurality of worlds'' is replaced by a single cost
  landscape.
  \item \textbf{Over Kripke.}  Actuality is not a matter of stipulation
  or rigid designation but is structurally determined by the cost
  geometry: $x = 1$ is the unique zero-cost configuration.
  \item \textbf{Over Williamson.}  While Williamson argues for
  necessitism on abductive grounds, RS \emph{derives} the unique
  necessary configuration from a functional equation.  The existence of
  necessities is proved, not postulated.
  \item \textbf{Over Plantinga.}  Maximal states of affairs are replaced
  by cost-geometric structure, giving modal concepts continuous,
  quantitative content via modal distance.
  \item \textbf{Over Stalnaker.}  Stalnaker's moderate realism about
  propositions as sets of possible worlds is refined: possible
  ``worlds'' are configurations with specific, computable $J$-costs, and
  the proposition-set structure inherits a metric from $J$.
\end{itemize}

\subsection{Necessitarianism: Spinoza and Leibniz}

The RS identification of actuality with necessity---the claim that the
actual world is the \emph{only} necessary one---is a form of
necessitarianism.  This has deep historical resonance.

Spinoza (\emph{Ethics}, Part~I) argued that everything follows
necessarily from the nature of God/Nature, so there is no genuine
contingency.  In RS, this is expressed geometrically: $x = 1$ is the
unique zero-defect state, and its existence is structurally forced by the
cost functional.

Leibniz asked ``Why is there something rather than nothing?'' and
answered: because God chose the best of all possible worlds.  RS gives
a non-theistic, geometric version:
\begin{equation}
  J(1) = 0 \quad\text{while}\quad J(x) \to \infty
  \;\text{as}\; x \to 0^+.
\end{equation}
``Something'' ($x = 1$) has zero cost; ``nothing'' ($x \to 0$) has
infinite cost.  Existence is cheaper than non-existence.  The actual
world exists not because it was chosen but because non-existence would
require infinite defect.

RS necessitarianism is, however, more constrained than Spinoza's.  RS
distinguishes sharply between the \emph{laws} (which are necessary:
there is a unique cost-zero configuration) and the
\emph{configurations with nonzero defect} (which are merely possible).
The merely possible configurations are real in the sense of having finite
cost, but they are not actual---they represent departures from the
necessary state.  This yields a compatibilist picture: the laws are
necessary, but the specific state of the world at any given tick is
contingent upon the dynamics.

\subsection{Counterfactuals}

The framework naturally accommodates counterfactual reasoning.  A
\emph{counterfactual configuration} relative to $c$ is any
$c' \in \mathrm{Poss}(c)$ with $c' \neq A(c)$ and
$J(c'.\,x) > J(A(c).\,x)$.  These are configurations that \emph{could
have} been reached (they are in the possibility set) but were not
actualized (they have higher cost than the actualization).

The ``closeness'' of a counterfactual to actuality is measured by the
cost difference $J(c'.\,x) - J(A(c).\,x)$: the smaller this value, the
``closer'' the counterfactual scenario is to what actually occurred.
This replaces Lewis's similarity metric on possible worlds with a
cost-geometric measure grounded in the functional $J$.

\subsection{Falsifiability}

The RS modal grounding is empirically constrained in three ways:
\begin{enumerate}
  \item \textbf{Uniqueness of necessity.}  If a physical system were
  found to have multiple zero-defect states ($J(x) = 0$ for
  $x \neq 1$), the uniqueness theorem would be falsified.  RS predicts
  that every physical observable, when expressed as a ratio to its
  natural unit, has $x = 1$ as its unique zero-cost value.
  \item \textbf{Positivity constraint.}  If a physical ratio $x$ were
  measured to be non-positive, the possibility/impossibility boundary
  would be falsified.  RS predicts all physical ratios are strictly
  positive.
  \item \textbf{Cost monotonicity.}  RS predicts that actualization
  decreases cost (Proposition~\ref{prop:cost-decrease}).  A physical
  process that systematically increased $J$-cost without external
  coupling would falsify the framework.
\end{enumerate}

%% ============================================================
\section{Conclusion}
\label{sec:conclusion}
%% ============================================================

Modal logic is grounded in the cost functional
$J(x) = \tfrac{1}{2}(x + x^{-1}) - 1$, the unique normalized solution
to the Recognition Composition Law.  The key results:
\begin{enumerate}
  \item Necessity is the unique $J$-minimizer: only $x = 1$ is
  necessary (Theorem~\ref{thm:unique-nec}).
  \item Possibility is positive ratio: every $x > 0$ is possible
  (Theorem~\ref{thm:possible}).
  \item Impossibility is the positivity boundary: $x \leq 0$
  (Theorem~\ref{thm:impossible}).
  \item Actuality is the $J$-minimum, not an indexical label
  (Theorem~\ref{thm:actual-unique}).
  \item The implication chain
  necessary $\Rightarrow$ actual $\Rightarrow$ possible holds
  (Corollary~\ref{cor:chain}).
  \item Every possible configuration is accessible to the necessary
  one via cost-decreasing paths (Theorem~\ref{thm:B}).
  \item Modal distance $d(x) = J(x)$ quantifies departure from
  actuality (Theorem~\ref{thm:distance}).
  \item Propositional operators $\Box$ and $\Diamond$ satisfy modal
  duality, axiom~K, and weak axiom~T
  (Propositions~\ref{prop:duality}--\ref{prop:T-weak}).
\end{enumerate}

This gives modal metaphysics physical content without ontological
inflation.  ``Necessary'' means $J$-forced.  ``Possible'' means
finite cost.  ``Actual'' means $J$-minimum.  Modal distance is
$J$-defect.

The framework entails a form of necessitarianism grounded in cost
geometry: the actual world is the unique necessary one, existence is
cheaper than non-existence, and the entire modal structure emerges from
a single functional equation.

%% ============================================================
\begin{thebibliography}{99}

\bibitem{Lewis1986}
D.~Lewis, \emph{On the Plurality of Worlds} (Blackwell, 1986).

\bibitem{Kripke1980}
S.~A.~Kripke, \emph{Naming and Necessity} (Harvard University Press,
1980).

\bibitem{Plantinga1974}
A.~Plantinga, \emph{The Nature of Necessity} (Clarendon Press, 1974).

\bibitem{Williamson2013}
T.~Williamson, \emph{Modal Logic as Metaphysics} (Oxford University
Press, 2013).

\bibitem{Spinoza1677}
B.~Spinoza, \emph{Ethics} (1677), translated by E.~Curley (Princeton
University Press, 1994).

\bibitem{Stalnaker2012}
R.~Stalnaker, \emph{Mere Possibilities: Metaphysical Foundations of
Modal Semantics} (Princeton University Press, 2012).

\bibitem{Divers2002}
J.~Divers, \emph{Possible Worlds} (Routledge, 2002).

\bibitem{Fine2005}
K.~Fine, \emph{Modality and Tense: Philosophical Papers} (Oxford
University Press, 2005).

\bibitem{Aczel1966}
J.~Acz\'el, \emph{Lectures on Functional Equations and Their Applications},
Academic Press, New York (1966).

\bibitem{RS_Axioms}
J.~Washburn, ``The Algebra of Reality: A Recognition Science Derivation
of Physical Law,'' \emph{Axioms} \textbf{15}(2), 90 (2026).
\texttt{doi:10.3390/axioms15020090}

\end{thebibliography}

\end{document}
